% !PRINTED 690
\weberpage{690}
%<<CONT
Einheitswurzeln bedeutet, und wenn wir also diese $w$ associirten Zahlen zu einem Complex zusammenfassen, so ist die Anzahl dieser Complexe gleich der Anzahl $T$ aller nicht associirten, durch $\mathfrak{a}$ theilbaren ganzen Zahlen in $\Omega$, deren absolute Norm kleiner als $t$ ist. Demnach ist die Anzahl der in $S'$ liegenden Gitterpunkte
\weq{14}{Z_t = w T.}

Für die Begrenzung des Gebietes $S$ erhält man aus (12) und (13):
\weq{15}{0 \leqq \xi_1 < 1, \ldots, 0 \leqq \xi_{\nu-1} < 1, \xi_\nu < 0,}
und nach (5) erhalten wir
\weq{16}{\mathrm{Lim}\, \frac{T}{t} = \frac{1}{w} \int\!\int \ldots \int dx_1\, dx_2 \ldots dx_n = \frac{1}{w} V,}
worin das $n$-fache Integral über das durch (15) bestimmte Gebiet $S$ auszudehnen ist.

\begin{center}
\S.\,190.

\gesperrt{Bestimmung des Volumens.}
\end{center}

Das Volumen $V$ lässt sich bestimmen nach der schon im \S.\,155 angeführten Transformationsformel für mehrfache Integrale, nach der, wenn $x_1, x_2, \ldots, x_n$ Functionen der Variablen $u_1, u_2, \ldots, u_n$ sind,
\weq{1}{\begin{aligned}
V &= \int\!\int \ldots \int dx_1\, dx_2 \ldots dx_n\\
&= \int\!\int \ldots \int \left( \Sigma \pm \frac{\partial x_1}{\partial u_1} \frac{\partial x_2}{\partial u_2} \cdots \frac{\partial x_n}{\partial u_n} \right) du_1\, du_2 \ldots du_n
\end{aligned}}
ist, wobei die Functionaldeterminante in dem nach den Variablen $u$ genommenen Integrale mit ihrem absoluten Werthe zu nehmen ist.

Diese Functionaldeterminante ist auch gleich dem reciproken Werthe der Determinante
\[ \Sigma \pm \frac{\partial u_1}{\partial x_1} \frac{\partial u_2}{\partial x_2} \cdots \frac{\partial u_n}{\partial x_n}. \]

Wir führen zunächst die Variablen $u$ durch eine lineare Substitution ein, und bemerken, dass unter den mit $\Omega$ conjugirten Körpern nach unserer Annahme $n - \nu$ conjugirt imaginäre Paare
%CONT>>
